% Contesto e Motivazioni
L'ecosistema blockchain di Algorand, pur offrendo elevate prestazioni e sicurezza nativa, sconta attualmente una severa carenza di strumenti automatizzati per l'analisi dei suoi Smart Contract, sviluppati tramite il linguaggio PyTeal. Parallelamente, i recenti Large Language Models (LLM) hanno dimostrato grandi potenzialità nell'analisi del codice, ma se applicati in modalità \textit{zero-shot} alla sicurezza software tendono a degenerare nella ricerca di correlazioni sintattiche superficiali, risultando proni ad allucinazioni e incapaci di dedurre la reale logica causale delle vulnerabilità. 


% Metodologia
Per superare queste limitazioni, questa tesi propone e valuta una pipeline di \textit{Vulnerability Detection} ibrida e \textit{LLM-Centric}, specificamente ottimizzata per l'ecosistema Algorand. Il framework analizza l'impatto di due metodologie avanzate: l'adattamento al dominio tramite \textbf{Fine-Tuning supervisionato basato su esempi Chain-of-Thought (CoT)} --- mirato a istruire il modello a generare un ragionamento analitico passo-passo --- e l'integrazione di contesto dinamico tramite \textbf{Retrieval-Augmented Generation (RAG)}. Per ovviare ai bias semantici del recupero basato su codice grezzo, il modulo RAG implementato sfrutta un'\textbf{astrazione semantica} ad alto livello del codice. Questa astrazione abilita un \textit{In-Context Learning} mirato: recuperando esempi di codice vulnerabile pertinenti e descrizioni di sicurezza per iniettarli direttamente nel \textit{prompt}, il sistema fornisce all'LLM riferimenti logici chiari per guidarne l'analisi.

% Risultati
I risultati sperimentali dimostrano l'efficacia dell'approccio proposto, evidenziando in particolare il netto incremento prestazionale apportato dalle metodologie rispetto agli approcci tradizionali. L'integrazione del RAG si è rivelata di grande impatto: agendo come meccanismo di validazione fattuale e abbattendo drasticamente i falsi positivi, ha garantito un miglioramento sostanziale dell'F1-Score e della Recall rispetto alle \textit{baseline} in \textit{zero-shot}. Questo ha permesso ai modelli di grossa taglia (come DeepSeek V3.2) di rasentare l'accuratezza ideale senza necessità di addestramento specifico. Al contempo, l'addestramento specifico ha ridotto sensibilmente la tendenza alle allucinazioni dei modelli di taglia media (come \texttt{QwQ-32B}), fornendo loro le competenze necessarie per condurre i task di audit in maniera molto più rigorosa e approfondita rispetto alla controparte base.

% Conclusioni
I risultati sperimentali evidenziano come la necessità del Fine-Tuning dipenda fortemente dalla dimensione del modello. Se nei modelli massivi l'\textit{In-Context Learning} tramite RAG è sufficiente per ottenere risultati ottimali,
per i modelli di taglia media l'addestramento specifico rimane determinante per condurre audit rigorosi in ambienti privati (\textit{on-premise}), colmando di fatto il distacco rispetto alle più costose soluzioni commerciali.

Il presente lavoro costituisce uno dei primi framework di sicurezza basati su IA ottimizzati per PyTeal. Al fine di garantire la totale riproducibilità scientifica, l'intero codice sorgente, i dataset generati e le configurazioni dei modelli sono stati resi pubblicamente disponibili con licenza open-source.
